\chapter{Discussion}
\label{chap:discussion}

\section{Architectural Findings}

The implementation and runtime traces show that modular planning can be inspected at boundaries that are absent from a prompt-to-action design. The planner does not own robot execution, speech output, or raw perception. It receives structured requests and returns structured plans or dialogue acts. Direct atomic execution was retained as an orchestrator control seam, but it was not treated as a matched experimental baseline because its capability scope differs from multi-step planner mediation.

The most useful architectural seam is the boundary between \code{planner_llm} and \code{nao_orchestrator}. The planner can propose a sequence, but the orchestrator validates and dispatches it. This makes it possible to test the planner with fake skills, physical skills, or unavailable skills without changing the planner prompt itself. It also means that execution feedback can be represented as a first-class input to replanning rather than as an unstructured error message.

\section{Planner Behaviour}

The planner design is strongest when the task can be represented as a sequence of available skill contracts. Atomic commands, object search, navigation, and report-generation tasks map naturally to the registry. Composite instructions are harder because they require the planner to preserve order, carry evidence between steps, and decide whether a later step should still run after an earlier failure.

The deterministic failure profiles showed that recovery quality depends on both feedback lineage and the availability of a semantically different valid plan. Navigation fail-once cases produced newer plan versions and truthful closure in the strongest historical sweep. Pick and blocked-delivery cases were less stable, and some later traces reached a safe stop without a complete post-terminal report. The planner can therefore react to feedback, but the evidence does not support a general recovery guarantee.

\section{Grounding and World State}

Grounding remains a central limitation. The architecture distinguishes compact prompt context, the detector-derived scene summary, and raw KB evidence because each has a different role. The chatbot needs a bounded world-state projection, the planner needs actionable entity identifiers and targets, and operators need traceable source evidence. The F13 semantic audit showed that presence in the projection is insufficient when target roles are not preserved through selection, planning, dispatch, and reporting.

The stable KB insertion probes separated fact availability from model use. Name, colour, support, and count queries passed in broad runs, while grouped post-effect queries exposed stale location relations after apparent delivery success. Action completion and state correctness must therefore be evaluated independently. This distinction also prevents detector variability from being misclassified as a planner or semantic-grounding failure.

\section{ROS4HRI Modularity}

The stack is partly robot-agnostic and partly NAO-specific. The planner request, planner output, execution feedback, and dialogue act contracts are not inherently tied to NAO. The same architecture could be reused with another ROS-compatible platform if an equivalent skill registry and orchestrator adapter were provided. By contrast, speech execution, replay motion, posture control, head motion, and look-at behaviour depend on NAO-specific packages or adapters.

This separation is a practical strength. It allows the thesis to evaluate the planner loop independently from the physical robot while still preserving a credible path to robot execution. The fake-skill layer is especially valuable because it provides deterministic execution evidence while the physical skills remain subject to hardware, safety, and perception variability.

\section{Validation Implications}

The difference between the F13 trajectory result and semantic audit is the main methodological finding. A case can contain speech ingress, planner admission, execution feedback, a terminal marker, and final speech while still acting on the wrong targets. Full user-turn injection remains necessary for route correctness, dialogue ownership, and exactly-once speech. Planner and action-level probes remain useful for isolating plan admission and skill guards. Neither is sufficient without an oracle that compares requested, selected, planned, executed, and reported members and then verifies the resulting KB state.

The runtime-review index is useful for operational triage, but it is secondary to explicit denominators and acceptance rules. The earlier 8.2 assessment and later 6.0 cap describe different evidence depths rather than a uniform performance decline. Reporting both makes the change in theorem strength visible instead of suppressing a safety finding to preserve a headline score.
